\section*{AI Use Statement}
Generative AI tools assisted topic and gap discovery, conceptual framing and hypothesis refinement, literature discovery and organization, methodology and experiment-design critique, code drafting and experiment orchestration, formulation and checking of formal statements and proof drafts, result analysis and interpretation, adversarial manuscript review, artifact documentation, and language editing. AI outputs were not treated as sources of evidence or as autonomous authorship. The authors manually checked citations against primary sources, executed and validated the code and locked outputs, inspected the proof obligations, reviewed all AI-assisted text, and take responsibility for the complete submission.

\section*{Ethics Statement}
The work audits public benchmark responses and does not use clinical records or interact with deployed patients. We disclose attack prerequisites and negative results to avoid overstating private-label or clinical risk. The intended use is to improve the integrity of evaluation registries and active selection protocols.

\section*{Reproducibility Statement}
The anonymous artifact records the natural-alias lock, preregistered targets, paired seeds, complete robustness and defense grids, official-selector provenance, and machine-readable outputs. Its manifest binds the learned-audit history through Step~105, including the IMDB held-out primary and the prospective Yelp \texttt{NO\_GO}. It includes learned weights, pre-outcome prediction locks, sealed-outcome bindings, raw 3,000-pair trajectories, literal status labels, and independent reconstructions. Offline mode verifies every bundled byte and anonymity. Full mode rebuilds the paper and reconstructs both final studies' trajectories, inference, and gates; optional raw-input mode re-infers their roots and aliases when pinned public base weights are available. Public source packages, base weights, and third-party repositories are not embedded, and historical validators are not all rerun transitively.
